\section{Preliminaries}
  \subsection{Problem Definition}\label{prelim:definition}
    Let us formally define the market generation problem. Given a
    financial time series of asset price movements
    $\mathbf{X} = (\mathbf{X}_1,\ldots,\mathbf{X}_I)^T = (\mathbf{X}_t)_{t=1}^T\in \mathbb R^{T\times I}$
    with $I$ individual securities spanning time $T$, let $p(\mathbf{X})$ denote
    its real distribution. The objective of the market generator is to
    synthesize a time series $\hat{\mathbf{X}} = (\hat{\mathbf{X}}_t)_{t=1}^T$
    such that its distribution $q(\hat{\mathbf{X}})$ approximates $p(\mathbf{X})$ while
    preserving key stylized facts, generating diverse samples, and
    improving downstream task training and evaluation. More importantly
    we wish to evaluate the generalization capabilites of $q(\hat{\mathbf{X}})$
    against unseen data $p(\mathbf{X}_{\mathrm{new}})$.

  \subsection{Scope of Project}
    \subsubsection{Evaluated Methods}
      Our benchmark covers models from classical parametric methods to
      modern deep-learning architectures. We selected models that meet
      one of three criteria: (1) proven industry adoption, (2) research
      community acceptance, and (3) recent methodological innovation.

      Classical parametric models rely on predefined statistical
      distributions and assumptions about the data-generating process.
      They remain common in financial institutions because of their
      interpretability and strong theoretical foundations.

      We also evaluate modern deep-learning approaches, especially
      models that learn the data distribution directly from
      observations rather than assuming a specific form. These include
      generative adversarial networks and variational autoencoders,
      which represent the current state of the art in synthetic market
      data generation.

      An overview of the market-generator classes is provided in
      §\ref{sec:model-overview}.

    \subsubsection{Datasets}
      The benchmark aims to generate realistic and diverse synthetic
      time series from daily financial asset prices over 15 years
      (2010--2025), collected from Yahoo Finance. To evaluate
      cross-asset behavior, we include 20 US large-cap stocks from
      several sectors and 5 popular US index-tracking ETFs (see table
      \ref{tab:asset-classification})

        \begin{table}[t]
        \centering
        \caption{Asset classification used in StonkBench.}
        \label{tab:asset-classification}
        \begin{tabular}{l l l l}
        \hline
        Ticker & Name & Type & Sector/Exposure \\
        \hline

        AAPL & Apple & Stock & Technology \\
        MSFT & Microsoft & Stock & Technology \\
        NVDA & NVIDIA & Stock & Technology \\
        AVGO & Broadcom & Stock & Technology \\
        JPM & JPMorgan Chase & Stock & Financials \\
        LLY & Eli Lilly & Stock & Healthcare \\
        UNH & UnitedHealth Group & Stock & Healthcare \\
        AMZN & Amazon & Stock & Consumer Disc. \\
        TSLA & Tesla & Stock & Consumer Disc. \\
        CAT & Caterpillar & Stock & Industrials \\
        UNP & Union Pacific & Stock & Industrials \\
        META & Meta Platforms & Stock & Communication \\
        NFLX & Netflix & Stock & Communication \\
        PG & Procter \& Gamble & Stock & Consumer Staples \\
        COST & Costco Wholesale & Stock & Consumer Staples \\
        XOM & Exxon Mobil & Stock & Energy \\
        CVX & Chevron & Stock & Energy \\
        NEE & NextEra Energy & Stock & Utilities \\
        PLD & Prologis & Stock & Real Estate \\
        LIN & Linde & Stock & Materials \\
        SPY & S\&P 500& ETF & S\&P 500 \\
        QQQ & Nasdaq-100& ETF & Nasdaq-100 \\
        XLF & Finance & ETF & U.S. Financials \\
        IWM & Russell 2000& ETF & Small Caps \\
        XLV & Health Care & ETF & U.S. Healthcare \\
        \hline
        \end{tabular}
    \end{table}

    \subsubsection{Scope of Evaluation Taxonomical Criteria}
      Our evaluation framework follows the taxonomy proposed by
      Stenger et al. \cite{Stenger24} and incorporates measures from
      multiple papers in the field.

      We integrate evaluation metrics from several sources while
      maintaining a structured taxonomy, providing broad coverage of
      the key aspects of evaluation and setting a standard for future
      market-generator research. A full description of the evaluation
      suite is provided in §\ref{subsec:stats-measure} and
      §\ref{sec:utility}.
